\documentclass[dvipdfmx,a4paper]{article}
\usepackage{amsmath,amssymb,amsthm}
\usepackage[utf8]{inputenc}
\usepackage[dvipdfmx]{geometry}
\geometry{left=25mm,right=25mm,top=25mm,bottom=25mm}
\usepackage[japanese]{babel}

\title{二重時空理論における内在モジュラリティによる\\フェルマーの最終定理の仮想的証明スケッチ}
\author{二重時空理論における思考実験}
\date{2026年1月}

\begin{document}

\selectlanguage{japanese}

\maketitle

\begin{abstract}
本稿は、二重時空理論（非標準的な枠組みで、16次元双四元数代数 $\mathbb{H} \oplus \mathbb{H} \cong \mathrm{Cl}(3,1)$ にミンコフスキー時空の2つの可換コピーを埋め込むもの）の代数的構造とフェルマーの最終定理（FLT）を結びつける\emph{仮想的かつ推測的な}考察を提示する。

核心的なアイデアは、コンパクトな双対時空セクターが楕円曲線のモジュラリティに類似した\emph{内在モジュラリティ}を提供することである。$a^n + b^n = c^n$ ($n>2$) の非自明解が存在すると仮定すると、ねじれ不整合が生じ、この内在モジュラリティに違反し、離散モデル内で矛盾に至る。

\textbf{重要な免責事項}：これはFLTの厳密な証明\emph{ではない}。FLTは1994--1995年にアンドリュー・ワイルズ（リチャード・テイラーとの共同）によりモジュラー形式とガロア表現を用いて厳密に証明されている。本稿は二重時空理論のアイデアと数論的構造を結びつける純粋に推測的な思考実験であり、数学的な証明としての有効性はない。
\end{abstract}

\section{背景：二重時空理論}

二重時空理論は、すべての質量を持つ粒子が双四元数代数 $\mathbb{H} \oplus \mathbb{H} \cong \mathrm{Cl}(3,1)$（計量 $(-,,+,+,+)$）に内在的にエンコードされたミンコフスキー時空の対を持つと仮定する。

\begin{itemize}
  \item \textbf{通常時空}：基底 $\{j, kI, kJ, kK\}$。ベクトル $X = ct\, j + x\, kI + y\, kJ + z\, kK$。
  \item \textbf{双対時空}：基底 $\{k, jI, jJ, jK\}$。双対写像 $X \mapsto Xi$（擬スカラー $i$ による右乗）で関連し、時間の矢を反転させる。
  \item 完全ローター：$R_\text{total} = R_\text{usual} R_\text{dual}$、ここで
  \[
  R_\text{usual} = \exp\!\left(\frac{\theta}{2} \hat{p}\right), \quad \hat{p}^2 = +1,
  \]
  \[
  R_\text{dual} = \exp\!\left(\frac{\phi}{2} \hat{q}\right), \quad \hat{q}^2 = -1.
  \]
  \item ねじれ不整合ローター：$\Omega = R_\text{usual}^\dagger R_\text{dual}$。
  \item 不変量：キリング形式スカラー $J = \frac{1}{16} B(\log \Omega, \log \Omega)$。
\end{itemize}

有効ローレンツ因子は
\[
\gamma_\text{eff} = \cosh\!\frac{\theta}{2} \cos\!\frac{\phi}{2} - \sinh\!\frac{\theta}{2} \sin\!\frac{\phi}{2}
= \Re[\cos z] + \Im[\cos z], \quad z = \frac{\phi}{2} + i \frac{\theta}{2}.
\]

離散モデルでは $\theta_k, \phi_k \in \frac{2\pi}{N} \mathbb{Z}$（大きな有限 $N$）とし、整数双四元数環の単位群を有限にする。

\section{モジュラー形式の概要}

モジュラー形式は複素上半平面 $\mathbb{H} = \{\tau \in \mathbb{C} \mid \Im \tau > 0\}$ 上の解析関数 $f(\tau)$ で、モジュラー群 $\mathrm{SL}(2,\mathbb{Z})$ の作用
\[
f\left( \frac{a\tau + b}{c\tau + d} \right) = (c\tau + d)^{k} f(\tau) \quad (a,b,c,d \in \mathbb{Z}, \, ad-bc=1)
\]
の下で重み $k$ の変換性を持つものである。特に重み $0$ のモジュラー関数は不変である。

代表例が $j$ 不変量で、すべてのモジュラー関数は $j(\tau)$ の有理関数として表現される。$j(\tau)$ は $q = \exp(2\pi i \tau)$ によるローラン展開
\[
j(\tau) = q^{-1} + 744 + 196884 q + 21493760 q^2 + \cdots
\]
を持ち、$\Im \tau \to +\infty$（$q \to 0$）で急速収束し、基本領域内で一価である。

ワイルズの証明では、有理数体上の半安定楕円曲線がすべてモジュラー形式から生じる（モジュラリティ定理）ことが鍵であり、非モジュラーなFrey曲線が矛盾を生む。

\section{双対時空のコンパクト性による内在モジュラリティ}

双対セクターはコンパクト（$\phi \mod 4\pi$）、通常セクターは非コンパクト（$\theta \in \mathbb{R}$）である。この非対称性は内在的なモジュラー様不変性を誘発する。

上半平面への写像を定義する：
\[
\tau = \frac{z}{\pi} = \frac{\phi}{2\pi} + i \frac{\theta}{2\pi}.
\]
$x = \phi/(2\pi)$（周期的）、$y = \theta/(2\pi)$（双曲的）とする。

双対ローターの周期性 $\phi \to \phi + 4\pi$ は $x \to x + 2$ に対応する。双四元数の代数的構造と合わせ、$\tau$ の $\mathrm{SL}(2,\mathbb{Z})$ に類似した部分群の下での変換は、ねじれ不整合を離散単位まで不変に保つ。

$q = \exp(2\pi i \tau) = \exp(i \phi - \theta)$ は、$\theta$ 大で急速収束し、$\gamma_\text{eff}$ の双曲・三角関数展開と類似の構造を示す。

したがって、本理論は\emph{内在モジュラリティ}を持つ：物理状態は $\tau$ のモジュラー不変配置に対応し、$j(\tau)$ に類似した不変量がねじれ不整合 $J$ として現れる。

\section{仮想的矛盾の仮定}

本思考実験では、フェルマーの最終定理の対偶を証明する形で議論を進めるため、$n > 2$ の整数に対し、$a^n + b^n = c^n$ を満たす原始非自明正整数解 $(a, b, c, n)$ が存在すると仮定する。ここで「原始」とは $\gcd(a, b, c) = 1$ であり、「非自明」とは $a, b, c > 0$ を意味する。この仮定の下で、二重時空理論の代数的構造内に矛盾が生じることを示す。

まず、通常時空セクターの単位ブーストバイベクトル $\hat{p} \in \mathrm{span}\{iI, iJ, iK\}$ を次のように構成する：
\[
\hat{p} = \frac{a^n \, iI + b^n \, iJ + c^n \, iK}{\sqrt{(a^n)^2 + (b^n)^2 + (c^n)^2}}.
\]
仮定された方程式 $a^n + b^n = c^n$ により、分母は厳密に
\[
\sqrt{(a^n)^2 + (b^n)^2 + (c^n)^2} = \sqrt{(c^n)^2 + (c^n)^2} = c^n \sqrt{2}
\]
となり、$\hat{p}^2 = +1$ が厳密に成立する。すなわち、$\hat{p}$ は通常時空の正規化されたブースト生成バイベクトルである。

次に、対応するラピディティ $\theta$ を
\[
\theta = 2 \cdot n \cdot \log\!\left( \frac{c}{\gcd(a,b,c)} \right) = 2 n \log c
\]
（原始性により $\gcd(a,b,c)=1$）と定義する。この選択は、$n$ 次高次性が通常時空の急速ブーストを異常スケーリングさせることを反映したものである。具体的には、$\theta / 2 \propto n \log c$ により、$n > 2$ の高次指数が非線形に虚部を増大させる。

上半平面への写像 $\tau = z / \pi$（$z = \phi/2 + i \theta/2$）の下で、対応する $\tau$ の虚部は
\[
y = \Im \tau = \frac{\theta}{2\pi} = \frac{n \log c}{\pi}
\]
となる。$n > 2$ かつ $c \gg 1$（非自明解の典型）により、$y \propto n$ は $n$ に比例して異常増大する。

この異常増大は、$q = \exp(2\pi i \tau) = \exp(i \phi - \theta)$ の収束を破壊的に加速し、ねじれ不整合ローター $\Omega = R_\text{usual}^\dagger R_\text{dual}$ の $q$ 展開類似で非収束項（異常極）を誘発する。特に、離散モデル $(\theta_k, \phi_k \in 2\pi \mathbb{Z}/N)$ では、$\hat{p}$ の方向が高次有理点 $(a^n : b^n : c^n)$ に固定されるため、有限 $N$ の格子点がこの異常スケーリングを収容できず、$\Omega$ が整数双四元数環の有限単位群外に逸脱する。

この構成は、双対ローター $R_\text{dual}$（コンパクトで有限離散位相）と組み合わせた際に、ねじれ不整合ローター $\Omega = R_\text{usual}^\dagger R_\text{dual}$ が高次 $n$ に依存した異常方向を取り、内在モジュラリティ（$\tau$ の保型不変性）を破る配置を強制する。以下節で、この非モジュラー挙動が離散モデル内で決定的矛盾に至ることを議論する。

\section{$n > 2$ における非モジュラー挙動}

$n=2$ の場合（ピタゴラス三つ組）をまず検討する。このとき、単位ブーストバイベクトル $\hat{p}$ の成分 $p_k \propto (a^2, b^2, c^2)$ は二次有理点であり、原始ピタゴラス三つ組の既知の生成公式（例：$a = m^2 - n^2$, $b = 2mn$, $c = m^2 + n^2$）により、有限個の離散スケールで単位球面上の有理点が生成される。これらの点は、上半平面写像 $\tau = \phi/(2\pi) + i \theta/(2\pi)$ の虚部 $y = \theta/(2\pi) \propto 2 \log c$ が有限範囲に収まり、モジュラー基本領域内に（離散同値変換まで）マッピングされる。双対ローターのコンパクト性により、ねじれ不整合ローター $\Omega = R_\text{usual}^\dagger R_\text{dual}$ は有限単位群内の元に留まり、$\log \Omega$ が定義され、キリング形式 $J = \frac{1}{16} B(\log \Omega, \log \Omega) \approx 0$ が達成可能である。これは特殊相対論の光円錐上を動く質量ゼロ粒子や、弱重力場での安定配置に対応し、内在モジュラリティを完全に保つ。

一方、$n > 2$ の場合、前節で構成した異常ブースト $\hat{p}$ とラピディティ $\theta \propto n \log c$ は決定的に異なる挙動を示す。高次指数 $n$ が単位球面上の有理点 $(a^n : b^n : c^n)$ を生成し、方向成分 $p_k$ は $n$ に依存した非線形スケーリングを持つ。これにより、対応する $q = \exp(2\pi i \tau) = \exp(i \phi - \theta)$ の絶対値 $|q| = e^{-\theta} \propto c^{-2n}$ が $n$ の増加とともに超指数的にゼロに収束する。

この超急速収束は、ねじれ不整合の $q$ 展開類似（$\gamma_\text{eff}$ の複素cos展開のフーリエアナロジー）で異常極を誘発する。具体的には：
\begin{itemize}
  \item $n=2$ では $q$ 展開の主要項が $q^{-1}$ 級で安定収束するが、
  \item $n > 2$ では高次 $n$ が $q^{-n}$ 級の非物理的発散項を生み、モジュラー関数 $j(\tau)$ のローラン展開（$q^{-1} + 744 + \cdots$）に類似した不変量が異常極を持ち、保型性を破る。
\end{itemize}

離散モデル（$\theta_k, \phi_k \in \frac{2\pi}{N} \mathbb{Z}$、固定大整数 $N$）では、この異常が決定的矛盾に至る。整数双四元数環の単位群は有限個の元しか持たず、$\Omega$ が取り得る値は有限集合に制限される。しかし $n > 2$ の高次有理点は、無限降下的な方向精緻化を強制し、$\Omega$ を単位群外に押しやる：
\begin{itemize}
  \item $\log \Omega$ が離散環内で未定義となるか、
  \item 無理やり近似しても $J$ が離散許容値外（無限大発散または非物理的負値爆発）を取り、トーション星の即時崩壊や特異点形成を誘発。
\end{itemize}

この非モジュラー挙動は、双対時空のコンパクト性（内在モジュラリティ）が強制する保型不変性に直接違反する。$n=2$ では許容される安定配置が、$n > 2$ の高次異常により原理的に不可能となる。

\section{思考実験の結論}

$n>2$ の非自明解の仮定は、離散モデル内の二重時空理論の内在モジュラー不変性と矛盾する配置に至る。

したがって、そのような解は存在し得ない。

\textbf{再強調}：これは有効な証明ではない。あくまで二重時空代数とモジュラー形式の構造的類似を示す推測的アナロジーである。FLTの本当の証明は20世紀の深い数論に依拠する。

\end{document}